
\subsubsection{6.1 Key Findings and Their Implications}

\textbf{SE and min\_logprob measure genuinely distinct uncertainty dimensions.} The Pearson |r|=0.026 is not merely a marginal pass of the 0.70 threshold --- it indicates near-orthogonality by any standard metric for variable collinearity. This gap from the prior first-token measurement (r=0.54--0.76) is mechanistically explained: min\_logprob captures late-sequence token uncertainty while first-token confidence captures only the initial prediction. The operational separation between SE (stochastic sampling, semantic aggregation) and min\_logprob (greedy decode, token-level minimum) manifests as near-zero empirical correlation.

This finding establishes a necessary precondition for ensemble benefit: it rules out the scenario where combining SE and min\_logprob provides no benefit because one is a redundant reparameterization of the other. However, necessity is not sufficiency --- and the partial $R^2$ null result (0.0005 at N=2500) shows that independence does not automatically yield linear ensemble gain. The signals are empirically near-orthogonal yet carry redundant correctness-predictive information in the linear logistic regression setting.

\textbf{Cross-model judge design suppresses circularity effectively.} The $\rho$(SE, judge) = $-0.026$ result validates the evaluation protocol: using Qwen-2.5-7B to evaluate Llama-3.1-8B outputs prevents shared model-specific NLI biases from creating circular agreement. This design choice, recommended by Santilli et al. (2025), is empirically validated in this setting.

\textbf{The partial $R^2$ null result: power characterization and gate disclosure.} The partial $R^{2}=0.0005$ at N=2500 (LRT p=0.442) was classified by the pre-registered pipeline gate as \texttt{EXPLORE\textbackslash \_N10} (gate\_pass: false). This outcome is disclosed transparently. At N=2500, estimated statistical power exceeds 0.99 for detecting the pre-registered effect size of partial $R^{2}\geq 0.02$ at $\alpha =0.05$. The gate's EXPLORE\_N10 flag reflects the pipeline's conservative joint marginal-uncertainty criterion, not a power inadequacy at the pre-registered threshold. We therefore characterize the partial $R^{2}=0.0005$ finding as a powered null for the pre-registered effect size: evidence of absence of a linear conditional effect at the pre-registered threshold on this benchmark, rather than absence of evidence due to underpowering. The dissociation between near-zero Pearson correlation (|r|=0.026) and near-zero conditional contribution (partial $R^{2}=0.0005)$ indicates that, while the signals are computed differently, they may capture overlapping correctness-predictive information in the linear logistic regression setting.

\subsubsection{6.2 Limitations}

\textbf{Ensemble AUROC not measured; linear null bounds linear gain.} Ensemble AUROC is not evaluated in this work. The partial $R^2$ null at N=2500 (gate=EXPLORE\_N10) indicates linear logistic regression ensembles will not achieve the pre-registered $\Delta AUROC$ target of $\geq 0.025$. Ensemble AUROC on TriviaQA dev and cross-model transfer to Qwen-2.5-7B on TruthfulQA are neither confirmed nor refuted; these are primary targets for follow-up work (h-e1-v2), which will also extend to N=10 stochastic samples per prompt as the gate recommends.

\textbf{Single seed and no bootstrapped confidence intervals.} N=2500 with a fixed seed (42) provides no estimate of result variance. The Pearson |r|=0.026 and Spearman $\rho=-0.026$ are point estimates. Bootstrapped 95\% CIs at N=2500 are planned for the h-e1-v2 analysis alongside nonlinear ensemble evaluation.

\textbf{Scope: short-answer factual QA with 7--8B open-weight models.} Results are validated on TriviaQA dev ($\leq 50$ token answers) with Llama-3.1-8B-Instruct. Generalization to long-form generation is explicitly out of scope --- Nguyen et al. (2025) document SE's weakness for multi-sentence answers.

\textbf{Unverified references.} The Gabriel (2026) paper (arXiv:2605.05166) and Santilli et al. (2025) (arXiv:2504.13677) were not found in Semantic Scholar at the time of writing; they are cited as preprints. The comparison finding from Gabriel (2026) (r=0.54--0.76 for first-token confidence) is used only to motivate the contrast with min\_logprob, not as a load-bearing empirical claim.

\subsubsection{6.3 Broader Impact}

This work contributes to the safety and reliability of LLM deployments by establishing the empirical independence of two lightweight UQ signals --- one requiring only greedy decode access (min\_logprob) and one requiring only generation API access (SE\_N5). The replicable evaluation methodology (cross-model LM-judge, circularity diagnostic, mechanism reality checks) provides a template for future UQ signal characterization studies. The pipeline uses publicly available open-weight models, ensuring reproducibility without proprietary API dependence.

